The North announced itself first as a smell, and only after that as a shape on the horizon.

Tyrion Lannister had noticed this on the voyage up from Dragonstone and had not yet decided what to make of it. The Sunset Sea south of the Neck smelled of salt and tar and, when the wind was contrary, of the seven hundred soldiers, forty horses, and two dragons currently occupying a fleet never designed to carry either of them. North of the Neck, the salt remained, but something else crept beneath it — pine and peat-smoke and a cold so clean it seemed to scour the nose of every other odor a man had grown used to breathing. He had said as much to Varys, three days back, and Varys had given him the small, unreadable smile he reserved for observations he considered beneath comment.

He was thinking of that smile now, watching the coastline resolve itself out of the grey morning — the long dark humps of hills, the white smudges that would become, in another hour or two, the towers of White Harbor — when Grey Worm came to stand beside him at the rail.

They had an understanding, he and this soldier, though neither had ever troubled to put a name to it. It was not friendship, exactly. Tyrion was not certain the Unsullied had been left the parts of themselves that friendship required, though he had come, over two years of proximity, to suspect this was a lie the world told about them more than a truth about what had actually been cut away. It was something more like the understanding between two men who have both, at different times, been made into instruments for other people’s purposes, and who recognize the particular ache of that in each other without needing to discuss it.

“You are thinking of her,” Tyrion said, because it was easier to say than any of the things he was actually thinking, and because it was true, and because Grey Worm’s eyes had not left the northern shore since it first came into view, in a way that had nothing to do with strategy.

“Missandei is below,” Grey Worm said. “With the Queen.”

“That isn’t what I asked.”

Grey Worm did not answer for a long moment. The wind pulled at the grey wisps of his hair, what little of it the years and the Masters had left him. When he spoke again his Common Tongue was as careful and deliberate as it always was, a man laying stones one at a time across a river he did not entirely trust to hold his weight.

“I have crossed water before,” he said. “The Master’s ships, when I was a boy, and no choosing where they took me. Now I cross water and I have chosen it. I did not think, when I was small, that this would ever be a thing I could say.”

“And now that you can say it?”

“Now I find I am afraid of what is on the other side of it. I did not used to be afraid of anything. Fear was cut from me with everything else, they said, when I was a boy in the fighting pits.” He glanced sideways, and there was something almost rueful in it, an expression so nearly human that Tyrion felt an unexpected tightness in his chest. “I do not think that was ever true. I think I was only very small, and there was nowhere to put the fear, so it went somewhere I could not find it. Now there is a place to put it.”

“And where is that?”

“Her,” Grey Worm said simply, and did not elaborate, and did not need to.

Tyrion looked back at the water rather than at his companion’s face, out of a courtesy he did not often extend to anyone, himself least of all. He thought of Tysha, briefly, unwillingly, the way he still did at odd unguarded moments some fifteen years on — a habit of grief he had never managed to break himself of, the mental equivalent of a missing tooth the tongue cannot leave alone. He thought that he understood, perhaps better than most men living, what it was to have found the one place fear could be put, and to have had it taken from him afterward with such surgical cruelty that he had spent a decade and a half telling himself it had never really been there in the first place.

He did not say any of this to Grey Worm. Some kindnesses were best delivered as silence.

“For what it is worth,” he said instead, “I have found, in a long and largely disgraceful life, that the things worth being afraid of losing are generally the only things worth having gained at all. I would not call that comforting. I would call it accurate, which is sometimes the best a man can do.”

Grey Worm considered this with the same grave attention he seemed to give everything, weighing it, turning it, deciding at last that it was a stone that would hold his weight. “That is very sad,” he said.

“Most true things are. It’s why men prefer the false ones.”

That, at least, won him something like the ghost of a smile — gone almost before it arrived, but real while it lasted. They stood together a while longer without speaking, watching White Harbor grow from smudge to shape to city, the seal-and-trident banners of House Manderly just visible now above the grey stone walls, and the fires of a garrison that had known, for weeks now, that a queen with dragons was coming, and had had entirely too much time to consider what that might mean for men who remembered dragons only from songs.

Grey Worm left him not long after, called below by some matter of the Unsullied companies and their disembarkation, and Tyrion had not been alone at the rail a quarter of an hour before Varys came to fill the space, soft-footed as ever, smelling faintly of the rosewater he affected even at sea, where Tyrion privately felt it did nothing but announce a eunuch’s presence to every gull within a mile.

“You have the look of a man doing sums he doesn’t like the answer to,” Varys said.

“I am doing several sums at once, actually, which is rather the trouble. One of them is soldiers. Another is Northern lords who have spent three hundred years hearing that dragons are extinct and are about to be disabused of the notion by a woman none of them have sworn to, arriving with an army twice the size of anything they raised against the Boltons. A third is what Sansa Stark makes of all this, which I confess worries me rather more than the first two combined.”

“Lady Stark is practical,” said Varys. “Practical people can be reasoned with.”

“Practical people can also be the most dangerous sort, when they’ve concluded that reason no longer serves them. I raised her, in a manner of speaking, for a year in King’s Landing, if ‘raised’ isn’t too generous a word for what my family did to that girl. She learned to smile at people who’d murdered everyone she loved. I’d not mistake that smile for warmth if I were you.”

Varys considered the coastline with the mild, patient expression he wore over nearly everything, the one Tyrion had long ago decided was either the most honest face in Westeros or the most thoroughly rehearsed. “And the Queen? Does she understand what she’s sailing into?”

“She understands there are Northern lords to win. I am considerably less certain she understands what it costs a people to bend the knee twice in three years to two different strangers, the second of whom arrives with three dragons and an army that speaks a language half of them have never heard. Robb Stark died for refusing a crown that wasn’t his to refuse twice. Their memory of kneeling is not a fond one.”

“You think she’ll fail to win them.”

“I think,” Tyrion said, and found, somewhat to his own surprise, that he meant it, “that she has never yet had to rule people who did not need her more than she needed them. Essos knelt to her because she broke their chains. The North does not need its chains broken. It needs to be convinced she’s worth the risk of trusting at all, which is an entirely different kind of persuading, and one I’m not certain either of us has taught her yet.”

Varys was quiet a moment, watching a gull wheel and drop toward the grey water. “You sound,” he said at last, mild as milk, “like a man auditing his own advice for the first time in some while.”

“I have found that a habit worth keeping, in my trade. The moment a man stops auditing his own advice is generally the moment it stops being worth anything.” Tyrion did not say the rest of what he was thinking — that he had given Daenerys a great deal of advice these past two years, and profited from most of it, and that he was no longer entirely sure whose interests he had been serving when he gave the best of it. That was a sum he preferred to do alone.

\sceneBreak

Jon found him not long before the ship’s bell rang for the turn into harbor — came up from belowdecks looking, Tyrion thought, like a man who had slept poorly and would go on sleeping poorly for some while yet, an expression Tyrion privately suspected had less to do with the dead marching south than with the more immediate matter of introducing a Targaryen queen to a sister who had spent her whole life being underestimated by people who mistook her courtesy for surrender.

“You’ve the look of a man who’d rather be facing the Night King than the docks ahead,” Tyrion said.

“I’d face him twice over before I face Sansa having to explain this to her.” Jon’s eyes were on White Harbor now too, the whole company of them fixed on that same grey shape as if staring hard enough might tell them something useful about what waited beyond it. “She trusted me to come back with an army. I’m not certain she pictured a Targaryen at the head of it.”

“Would it help if I told you Northern lords have survived stranger alliances?”

“It would help if it were true.”

“It is true. It simply won’t feel true to them until it’s worked, which is the trouble with every alliance ever made — they only look inevitable afterward.” Tyrion glanced sideways at him. “For what it’s worth, Lord Snow, I’ve watched a great many men walk into rooms expecting judgment and receive it. You have one advantage most of them lacked.”

“Which is?”

“You’re not lying to anyone in that room, including yourself. That’s rarer than you’d think, at the head of an army. Most men arrive already halfway convinced of their own performance. You look like a man still deciding whether he believes any of this will work, which is, perversely, the most convincing thing you have to offer them.” Tyrion allowed himself a small, tired smile. “Try not to lose it before we land.”

Jon didn’t answer that directly. He stood a while longer at the rail, hood down despite the cold, watching the towers of White Harbor sharpen out of the grey the way a wound sharpens as the numbness wears off it, and when he finally spoke again it was quieter, meant, Tyrion suspected, as much for himself as for the Hand beside him.

“I keep thinking of the Wall,” he said. “Whether it’s held. Whether there’s anyone left up there to tell us if it hasn’t.”

Tyrion had no comfort to offer that was not also a lie, and so, for once in a long practiced life of offering comfortable lies to powerful people, he offered none at all, and let the silence stand between them until the ship’s bell rang.

\sceneBreak

Below decks, in the queen’s cabin, Tyrion did not go, though he was fairly sure of what he would have found if he had. He had passed near the door twice in the last week at hours when a Hand of the Queen might reasonably be expected to require his queen’s attention, and twice he had heard, through wood thinned by salt and years, the low murmur of the King in the North’s voice and the quieter, rarer sound of Daenerys Targaryen laughing — not the practiced laugh she gave in audience, the one calibrated for effect, but something smaller and more startled, as though it kept catching her by surprise that she was still capable of it.

He was glad of it, and he was afraid of it, in roughly equal measure, and he suspected — though he had said this to no one, not even Varys, not even himself in so many words — that those two feelings were going to spend a great deal of the coming year refusing to resolve into just one.

A king who had bent the knee to a queen he was, by every report reaching Tyrion’s ears, quietly falling in love with, was a king whose knee might yet prove less bent than it looked, should the two feelings ever come into conflict. Tyrion had built enough of his life on reading the small cracks beneath grand gestures to know better than to relax simply because the ceremony had gone smoothly. Ceremonies always went smoothly. It was the years afterward that tested whether they had meant anything.

And there was the matter — delicate, unspoken, not yet a certainty even to the two people it most concerned, so far as Tyrion could judge — of what else two people confined to close quarters on a long voyage north might have made, besides an alliance.

He did not know. He suspected. He had learned, over a life spent watching rooms for a living, not to mistake suspicion for knowledge. But he found himself, watching White Harbor’s towers draw close under a sky the color of an old bruise — a color, he would learn only that evening, that stretched unbroken all the way north to a Wall that no longer stood — with a stomach that had survived a great many wars, a trial, a wildfire, and his own father’s crossbow, and was fairly sure, all the same, that it was not going to like whatever had broken that sky.